\documentclass[12pt,a4paper]{ctexart}
\usepackage{amsmath,amssymb}
\usepackage{geometry}
\usepackage{listings}
\usepackage{xcolor}
\usepackage{hyperref}
\hypersetup{
	colorlinks=true,
	linkcolor={blue!60!black},
	urlcolor={blue!60!black},
	citecolor={blue!60!black},
	pdfborder={0 0 0},
}
\usepackage{tabularx}
\usepackage{fancyhdr}
\usepackage{tikz}
\usepackage{lastpage}
\usepackage{enumitem}
\usepackage{array}
\usepackage{booktabs}
\usepackage{longtable}
\geometry{left=2.5cm,right=2.5cm,top=2.5cm,bottom=2.5cm}
\definecolor{codebg}{RGB}{245,245,245}
\definecolor{codeframe}{RGB}{200,200,200}
\definecolor{commentcolor}{RGB}{0,128,0}
\definecolor{keywordcolor}{RGB}{127,0,85}
\definecolor{stringcolor}{RGB}{42,0,255}
\definecolor{accent}{HTML}{4F46E5}
\definecolor{accent2}{HTML}{7C3AED}
\definecolor{mid}{HTML}{6B7280}
\lstset{
	language=C++,
	basicstyle=\small\ttfamily,
	backgroundcolor=\color{codebg},
	frame=single,
	rulecolor=\color{codeframe},
	breaklines=true,
	showstringspaces=false,
	commentstyle=\color{commentcolor},
	keywordstyle=\color{keywordcolor}\bfseries,
	stringstyle=\color{stringcolor},
	numbers=left,
	numberstyle=\tiny\color{gray},
	tabsize=4,
	captionpos=b,
	extendedchars=true,
	columns=flexible,
}
\pagestyle{fancy}
\fancyhf{}
\fancyhead[L]{\leftmark}
\fancyhead[R]{动态规划算法专题}
\fancyfoot[C]{\thepage\ /\ \pageref*{LastPage}}
\renewcommand{\headrulewidth}{0.4pt}
\renewcommand{\footrulewidth}{0pt}
\newcommand{\infobox}[1]{%
	\par\noindent
	\fcolorbox{accent!40}{accent!5}{%
		\begin{minipage}{\dimexpr\textwidth-2\fboxsep-2\fboxrule\relax}
			#1
		\end{minipage}%
	}\par
}
\title{动态规划算法专题}
\author{算法学习笔记}
\date{\today}
\begin{document}
	\maketitle
	\tableofcontents
	\newpage
	\section{线性DP}
	\subsection{最长上升子序列（LIS）}
	\infobox{解决问题：给定一个序列，求最长的严格递增子序列的长度。适用于股票走势分析、基因序列比对、导弹拦截系统等场景。}
	\subsubsection*{题目描述}
	给定一个长度为 $n$ 的数组 $a$，找出其中最长的严格递增子序列的长度。子序列不要求连续，但顺序必须与原数组一致。
	\subsubsection*{输入示例}
	\begin{verbatim}
		8
		3 1 2 6 4 5 10 7
	\end{verbatim}
	\subsubsection*{输出示例}
	5
	\subsubsection*{输出解释}
	最长上升子序列为：1, 2, 4, 5, 10（长度为5）
	\begin{lstlisting}
		#include <bits/stdc++.h>
		using namespace std;
		int main() {
			int n;  // n: 数组长度
			cin >> n;
			vector<int> a(n);  // a: 存储原始数组的向量
			for (int i = 0; i < n; i++) {
				cin >> a[i];  // 循环读入n个整数到数组a中
			}
			// dp[i]: 以a[i]结尾的最长上升子序列长度，初始化为1（子序列只包含自身）
			vector<int> dp(n, 1);
			int ans = 0;  // ans: 全局最长上升子序列的长度
			// 核心原理：对于每个位置i，遍历它前面所有位置j，如果a[j] < a[i]，
			// 则a[i]可以接在以a[j]结尾的序列后面，形成更长的上升子序列
			for (int i = 0; i < n; i++) {           // i: 当前考虑的子序列结尾元素下标
				for (int j = 0; j < i; j++) {       // j: 遍历i之前的所有元素下标
					if (a[j] < a[i]) {              // 判断a[j]是否小于a[i]，保证严格递增
						dp[i] = max(dp[i], dp[j] + 1);  // 状态转移：取最大值，接在j后长度+1
					}
				}
				ans = max(ans, dp[i]);  // 更新全局最长长度
			}
			cout << ans << "\n";  // 输出最终结果
			return 0;
		}
	\end{lstlisting}
	\subsection{最长公共子序列（LCS）}
	\infobox{解决问题：给定两个字符串/序列，求它们最长的公共子序列的长度。适用于版本对比、基因序列相似度分析、文件差异比较等场景。}
	\subsubsection*{题目描述}
	给定两个字符串 $s$ 和 $t$，找出它们的最长公共子序列的长度。公共子序列是指在两个字符串中都按顺序出现（不一定连续）的字符序列。
	\subsubsection*{输入示例}
	\begin{verbatim}
		abcde
		ace
	\end{verbatim}
	\subsubsection*{输出示例}
	3
	\subsubsection*{输出解释}
	最长公共子序列为 "ace"，长度为3
	\begin{lstlisting}
		#include <bits/stdc++.h>
		using namespace std;
		int main() {
			string s, t;  // s, t: 两个输入的字符串
			cin >> s >> t;
			int n = s.size(), m = t.size();  // n: 字符串s的长度, m: 字符串t的长度
			// dp[i][j]: s的前i个字符和t的前j个字符的最长公共子序列长度
			// i和j的取值范围都是从0到n/m，0表示空串
			vector<vector<int>> dp(n + 1, vector<int>(m + 1, 0));
			// 核心原理：遍历两个字符串的所有前缀组合
			// 若当前字符相等，则LCS长度在左上角基础上+1
			// 若不相等，则取左边或上边的最大值
			for (int i = 1; i <= n; i++) {      // i: s字符串的当前处理前缀长度
				for (int j = 1; j <= m; j++) {  // j: t字符串的当前处理前缀长度
					if (s[i - 1] == t[j - 1]) {  // 当前两个字符相等（注意字符串索引从0开始）
						// 相等时，LCS长度 = 去掉这两个字符后的LCS长度 + 1
						dp[i][j] = dp[i - 1][j - 1] + 1;
					} else {
						// 不相等时，LCS长度 = max(去掉s当前字符, 去掉t当前字符)
						dp[i][j] = max(dp[i - 1][j], dp[i][j - 1]);
					}
				}
			}
			cout << dp[n][m] << "\n";  // dp[n][m]即为两个完整字符串的LCS长度
			return 0;
		}
	\end{lstlisting}
	\subsection{最大子数组和（Kadane算法）}
	\infobox{解决问题：给定一个整数数组（可能包含负数），求连续子数组的最大和。适用于股票最大收益、最优区间选择、信号处理等场景。}
	\subsubsection*{题目描述}
	给定一个整数数组 $nums$，请找出一个具有最大和的连续子数组（子数组最少包含一个元素），返回其最大和。
	\subsubsection*{输入示例}
	\begin{verbatim}
		-2 1 -3 4 -1 2 1 -5 4
	\end{verbatim}
	\subsubsection*{输出示例}
	6
	\subsubsection*{输出解释}
	连续子数组 [4, -1, 2, 1] 的和最大，为6
	\begin{lstlisting}
		#include <bits/stdc++.h>
		using namespace std;
		int main() {
			int n;  // n: 数组长度
			cin >> n;
			vector<int> a(n);  // a: 存储原始整数的向量
			for (int i = 0; i < n; i++) {
				cin >> a[i];
			}
			int cur = a[0];   // cur: 以当前位置结尾的最大子数组和（当前最优）
			int ans = a[0];   // ans: 全局最大子数组和
			// 核心原理：Kadane算法，遍历数组时维护两个值
			// 对于每个元素，要么将其加入之前的子数组，要么以它作为新子数组的起点
			// cur = max(当前元素, 当前元素 + 之前的cur)
			for (int i = 1; i < n; i++) {  // i: 从第二个元素开始遍历
				// 状态转移：决定是开启新子数组(a[i])，还是延续旧子数组(cur + a[i])
				cur = max(a[i], cur + a[i]);
				ans = max(ans, cur);  // 用当前子数组和更新全局最大值
			}
			cout << ans << "\n";
			return 0;
		}
	\end{lstlisting}
	\subsection{最长回文子串}
	\infobox{解决问题：给定一个字符串，找出其中最长的回文子串。回文是指正着读和反着读都相同的字符串。适用于文本处理、DNA序列分析等场景。}
	\subsubsection*{题目描述}
	给定一个字符串 $s$，找出其中最长的回文子串。如果有多个，返回任意一个。
	\subsubsection*{输入示例}
	\begin{verbatim}
		babad
	\end{verbatim}
	\subsubsection*{输出示例}
	aba 或 bab
	\begin{lstlisting}
		#include <bits/stdc++.h>
		using namespace std;
		int main() {
			string s;  // s: 输入的原始字符串
			cin >> s;
			int n = s.size();  // n: 字符串的长度
			// dp[i][j]: 子串s[i..j]是否为回文串，true表示是回文
			vector<vector<bool>> dp(n, vector<bool>(n, false));
			int start = 0, maxlen = 1;  // start: 最长回文子串的起始索引, maxlen: 最长回文子串的长度
			// 核心原理：回文串去掉两端字符后仍是回文串
			// 先初始化所有长度为1和2的子串，再逐步扩展长度
			for (int i = 0; i < n; i++) {
				dp[i][i] = true;  // 所有单个字符都是回文串
			}
			for (int i = 0; i < n - 1; i++) {
				if (s[i] == s[i + 1]) {  // 相邻两个字符相等
					dp[i][i + 1] = true;  // 长度为2的子串是回文
					start = i;      // 更新最长回文起始位置
					maxlen = 2;     // 更新最长回文长度为2
				}
			}
			for (int len = 3; len <= n; len++) {          // len: 当前检查的子串长度，从3开始
				for (int i = 0; i + len - 1 < n; i++) {   // i: 子串的左边界索引
					int j = i + len - 1;                  // j: 子串的右边界索引
					// 当两端字符相等，且去掉两端后的内部子串是回文时，当前子串为回文
					if (s[i] == s[j] && dp[i + 1][j - 1]) {
						dp[i][j] = true;  // 标记[i, j]区间为回文串
						start = i;        // 更新最长回文串的起始位置
						maxlen = len;     // 更新最长回文串的长度
					}
				}
			}
			cout << s.substr(start, maxlen) << "\n";  // 截取并输出最长回文子串
			return 0;
		}
	\end{lstlisting}
	\section{背包DP}
	\subsection{0-1背包}
	\infobox{解决问题：有n个物品，每个物品只能选一次，在容量为W的背包中装入价值最大的物品组合。适用于资源分配、投资组合、货物装载等场景。}
	\subsubsection*{题目描述}
	有 $n$ 件物品和一个容量为 $W$ 的背包。第 $i$ 件物品的重量为 $w_i$，价值为 $v_i$。每件物品只能选择一次，求在不超过背包容量的前提下，能装入的最大总价值。
	\subsubsection*{输入示例}
	\begin{verbatim}
		4 5
		1 2
		2 4
		3 4
		4 5
	\end{verbatim}
	\subsubsection*{输出示例}
	8
	\subsubsection*{输出解释}
	选择物品1、2、3（重量1+2+3=6超重），选择物品1、2、4（重量1+2+4=7超重），最优是选择物品2和3（重量2+3=5，价值4+4=8）
	\begin{lstlisting}
		#include <bits/stdc++.h>
		using namespace std;
		int main() {
			int n, W;  // n: 物品的总数量, W: 背包的总容量
			cin >> n >> W;
			vector<int> w(n), v(n);  // w[i]: 第i个物品的重量, v[i]: 第i个物品的价值
			for (int i = 0; i < n; i++) {
				cin >> w[i] >> v[i];
			}
			// dp[j]: 背包容量为j时能装入的最大总价值，j的范围是0到W
			vector<int> dp(W + 1, 0);
			// 核心原理：对于每个物品，逆序遍历容量，确保每个物品只被选一次
			// 状态转移方程：dp[j] = max(不选当前物品, 选当前物品)
			for (int i = 0; i < n; i++) {          // i: 当前正在处理的物品的索引
				// 逆序遍历容量，防止同一个物品被重复计算（01背包特性）
				for (int j = W; j >= w[i]; j--) {   // j: 当前考虑的背包容量
					// dp[j]：不选第i个物品；dp[j-w[i]]+v[i]：选第i个物品
					dp[j] = max(dp[j], dp[j - w[i]] + v[i]);
				}
			}
			cout << dp[W] << "\n";  // 容量为W时的最大价值即为答案
			return 0;
		}
	\end{lstlisting}
	\subsection{完全背包}
	\infobox{解决问题：有n种物品，每种物品无限数量，在容量为W的背包中装入价值最大的物品组合。适用于零钱兑换、原材料无限供应等场景。}
	\subsubsection*{题目描述}
	有 $n$ 种物品和一个容量为 $W$ 的背包。第 $i$ 种物品的重量为 $w_i$，价值为 $v_i$，每种物品的数量是无限的。求在不超过背包容量的前提下，能装入的最大总价值。
	\subsubsection*{输入示例}
	\begin{verbatim}
		3 7
		2 3
		3 4
		4 5
	\end{verbatim}
	\subsubsection*{输出示例}
	10
	\subsubsection*{输出解释}
	选择2个物品1（重量2+2=4，价值3+3=6）加1个物品2（重量+3=7，价值+4=10）
	\begin{lstlisting}
		#include <bits/stdc++.h>
		using namespace std;
		int main() {
			int n, W;  // n: 物品种类的数量, W: 背包的总容量
			cin >> n >> W;
			vector<int> w(n), v(n);  // w[i]: 第i种物品的重量, v[i]: 第i种物品的价值
			for (int i = 0; i < n; i++) {
				cin >> w[i] >> v[i];
			}
			// dp[j]: 背包容量为j时能装入的最大总价值
			vector<int> dp(W + 1, 0);
			// 核心原理：对于每种物品，正序遍历容量，实现无限次选取
			// 因为正序遍历时，dp[j-w[i]]可能已经包含了当前物品，所以可以无限取
			for (int i = 0; i < n; i++) {           // i: 当前正在处理的物品种类的索引
				// 正序遍历容量，允许同一物品被多次选取
				for (int j = w[i]; j <= W; j++) {    // j: 当前考虑的背包容量
					// 与01背包的唯一区别就是内层循环是正序
					dp[j] = max(dp[j], dp[j - w[i]] + v[i]);
				}
			}
			cout << dp[W] << "\n";
			return 0;
		}
	\end{lstlisting}
	\subsection{多重背包（二进制优化）}
	\infobox{解决问题：有n种物品，每种物品有有限数量，在容量为W的背包中装入价值最大的物品组合。适用于库存有限、配额分配等场景。}
	\subsubsection*{题目描述}
	有 $n$ 种物品和一个容量为 $W$ 的背包。第 $i$ 种物品的重量为 $w_i$，价值为 $v_i$，数量为 $c_i$ 个。求在不超过背包容量的前提下，能装入的最大总价值。
	\subsubsection*{输入示例}
	\begin{verbatim}
		3 10
		2 3 3
		3 4 2
		4 5 2
	\end{verbatim}
	\subsubsection*{输出示例}
	18
	\subsubsection*{输出解释}
	物品1取3个(6), 物品2取1个(3), 物品3取1个(4), 价值=3*3+4*1+5*1=18
	\begin{lstlisting}
		#include <bits/stdc++.h>
		using namespace std;
		int main() {
			int n, W;  // n: 物品种类的数量, W: 背包的总容量
			cin >> n >> W;
			vector<int> w(n), v(n), c(n);  // w[i]: 第i种物品的重量, v[i]: 价值, c[i]: 该物品的最大数量
			for (int i = 0; i < n; i++) {
				cin >> w[i] >> v[i] >> c[i];
			}
			vector<int> dp(W + 1, 0);  // dp[j]: 容量为j的背包能装的最大价值
			// 核心原理：二进制优化，将每种物品按2的幂次分解成多个新的物品组
			// 这样将多重背包问题转化为01背包问题，时间复杂度降至O(W * Σlog(c[i]))
			for (int i = 0; i < n; i++) {   // i: 当前处理的物品种类索引
				int k = 1;          // k: 当前拆分的数量，从2^0=1开始
				int remain = c[i];  // remain: 当前剩余未拆分的物品数量
				while (remain > 0) {
					int take = min(k, remain);   // take: 当前组要取多少个物品
					int weight = w[i] * take;    // weight: 当前组的重量总和
					int value = v[i] * take;     // value: 当前组的价值总和
					// 将拆分出的这一组视为一个新的物品，用01背包的方式处理
					for (int j = W; j >= weight; j--) {
						dp[j] = max(dp[j], dp[j - weight] + value);
					}
					remain -= take;   // 从剩余数量中减去已取走的
					k <<= 1;          // k = k * 2，继续下一个2的幂次
				}
			}
			cout << dp[W] << "\n";
			return 0;
		}
	\end{lstlisting}
	\subsection{分组背包}
	\infobox{解决问题：物品分成若干组，每组中最多只能选一个物品，在容量为W的背包中装入价值最大的物品组合。适用于班级选课、产品线选择等场景。}
	\subsubsection*{题目描述}
	有 $n$ 组物品和一个容量为 $W$ 的背包。每组中有若干物品，每组最多只能选一个物品。第 $i$ 组第 $k$ 个物品的重量为 $w_{ik}$，价值为 $v_{ik}$。求能装入的最大总价值。
	\subsubsection*{输入示例}
	\begin{verbatim}
		3 5
		2 1 2 2 3
		2 2 3 3 4
		1 4 5
	\end{verbatim}
	\subsubsection*{输入解释}
	第1组有2个物品:(1,2)和(2,3)；第2组有2个物品:(2,3)和(3,4)；第3组有1个物品:(4,5)
	\begin{lstlisting}
		#include <bits/stdc++.h>
		using namespace std;
		int main() {
			int n, W;  // n: 物品组数, W: 背包总容量
			cin >> n >> W;
			vector<vector<int>> group_w, group_v;  // group_w[i]: 第i组所有物品的重量, group_v[i]: 第i组所有物品的价值
			for (int i = 0; i < n; i++) {   // i: 当前正在读取的组索引
				int cnt;  // cnt: 当前组内的物品数量
				cin >> cnt;
				vector<int> w(cnt), v(cnt);  // 临时存储当前组的重量和价值
				for (int j = 0; j < cnt; j++) {
					cin >> w[j] >> v[j];
				}
				group_w.push_back(w);
				group_v.push_back(v);
			}
			// dp[j]: 容量为j的背包能装的最大价值
			vector<int> dp(W + 1, 0);
			// 核心原理：先遍历组，再逆序容量，最后遍历组内物品
			// 逆序容量确保每组最多只选一个物品
			for (int i = 0; i < n; i++) {   // i: 当前正在处理的组的索引
				// 逆序遍历容量，防止同组物品被重复选择
				for (int j = W; j >= 0; j--) {   // j: 当前考虑的背包容量
					// 在当前组内遍历所有物品，尝试放入背包
					for (int k = 0; k < group_w[i].size(); k++) {  // k: 组内物品的索引
						if (j >= group_w[i][k]) {
							dp[j] = max(dp[j], dp[j - group_w[i][k]] + group_v[i][k]);
						}
					}
				}
			}
			cout << dp[W] << "\n";
			return 0;
		}
	\end{lstlisting}
	\section{区间DP}
	\subsection{石子合并}
	\infobox{解决问题：给定一排石子，每次只能合并相邻两堆，花费为两堆石子数之和，求合并成一堆的最小总花费。适用于最优括号匹配、矩阵链乘等场景。}
	\subsubsection*{题目描述}
	有 $n$ 堆石子排成一排，每堆石子有数量 $a_i$。每次可以合并相邻的两堆，花费为两堆石子数量之和。经过 $n-1$ 次合并后成为一堆，求最小总花费。
	\subsubsection*{输入示例}
	\begin{verbatim}
		4
		1 3 5 2
	\end{verbatim}
	\subsubsection*{输出示例}
	22
	\subsubsection*{输出解释}
	合并顺序：1,3合并(4)→现在为4,5,2→4,5合并(9)→9,2合并(11)，总花费=4+9+11=24；
	最优：1,3合并(4)→4,5,2→5,2合并(7)→4,7合并(11)，花费=4+7+11=22
	\begin{lstlisting}
		#include <bits/stdc++.h>
		using namespace std;
		int main() {
			int n;  // n: 石子堆的数量
			cin >> n;
			vector<int> a(n + 1);           // a[i]: 第i堆石子的数量，索引从1开始
			vector<int> prefix(n + 1, 0);   // prefix[i]: 前i堆石子的总重量（前缀和）
			for (int i = 1; i <= n; i++) {
				cin >> a[i];
				prefix[i] = prefix[i - 1] + a[i];  // 计算前缀和
			}
			// dp[i][j]: 将区间[i, j]内的石子合并成一堆的最小花费
			vector<vector<int>> dp(n + 2, vector<int>(n + 2, 0));
			// 核心原理：枚举区间长度，再枚举左端点，最后枚举区间内的分割点
			// 状态转移：dp[i][j] = min(dp[i][k] + dp[k+1][j] + sum(i,j))
			for (int len = 2; len <= n; len++) {        // len: 当前枚举的区间长度
				for (int i = 1; i + len - 1 <= n; i++) {  // i: 区间左端点
					int j = i + len - 1;                  // j: 区间右端点
					dp[i][j] = 1e9;  // 初始化为一个很大的值，用于求最小值
					for (int k = i; k < j; k++) {       // k: 分割点，将区间分成[i,k]和[k+1,j]
						// 合并总花费 = 左区间最小花费 + 右区间最小花费 + 当前合并花费(区间和)
						int cost = dp[i][k] + dp[k + 1][j] + (prefix[j] - prefix[i - 1]);
						dp[i][j] = min(dp[i][j], cost);
					}
				}
			}
			cout << dp[1][n] << "\n";  // 输出合并整个区间[1,n]的最小花费
			return 0;
		}
	\end{lstlisting}
	\subsection{括号匹配}
	\infobox{解决问题：给定一个括号序列，求最长有效括号子串的长度或最少添加几个括号使其合法。适用于代码语法检查、表达式求值等场景。}
	\subsubsection*{题目描述}
	给定一个只包含 '(' 和 ')' 的字符串，找出最长的有效（格式正确且连续）括号子串的长度。
	\subsubsection*{输入示例}
	\begin{verbatim}
		(()())
	\end{verbatim}
	\subsubsection*{输出示例}
	6
	\begin{lstlisting}
		#include <bits/stdc++.h>
		using namespace std;
		int main() {
			string s;  // s: 输入的括号字符串
			cin >> s;
			int n = s.size();  // n: 字符串的长度
			// dp[i]: 以第i个字符结尾的最长有效括号子串的长度
			vector<int> dp(n, 0);
			int ans = 0;  // ans: 全局最长有效括号长度
			// 核心原理：只处理')'，因为'('无法形成有效括号结尾
			// 分两种情况：'...()' 或 '...))...'，利用dp数组向前跳跃匹配
			for (int i = 1; i < n; i++) {  // i: 当前处理的字符索引，从1开始（0位置无法形成有效括号）
				if (s[i] == ')') {  // 只有右括号才可能形成有效括号结尾
					if (s[i - 1] == '(') {
						// 情况1: 形成 "...()" 模式
						// 有效长度 = 当前匹配的2 + i-2位置的有效长度
						dp[i] = (i >= 2 ? dp[i - 2] : 0) + 2;
					} else if (i - dp[i - 1] - 1 >= 0 && s[i - dp[i - 1] - 1] == '(') {
						// 情况2: 形成 "...))...)" 模式，需要找到匹配的左括号
						// i - dp[i-1] - 1 是可能匹配的左括号位置
						// 有效长度 = 内部有效长度 + 2 + 左括号前的有效长度
						dp[i] = dp[i - 1] + 2;
						if (i - dp[i - 1] - 2 >= 0) {
							dp[i] += dp[i - dp[i - 1] - 2];
						}
					}
					ans = max(ans, dp[i]);  // 更新全局最长长度
				}
			}
			cout << ans << "\n";
			return 0;
		}
	\end{lstlisting}
	\section{树形DP}
	\subsection{树的最大独立集（没有上司的舞会）}
	\infobox{解决问题：在树形结构中，选择不相邻的节点使得权值和最大。适用于公司年会邀请、选课系统、家庭聚会等场景。}
	\subsubsection*{题目描述}
	某公司有 $n$ 个员工，每个员工有一个快乐值。员工之间的上下级关系构成一棵树。如果邀请了某个员工，就不能邀请他的直接上级。求能获得的最大快乐值总和。
	\subsubsection*{输入示例}
	\begin{verbatim}
		7
		1 2 3 4 5 6 7
		1 2
		1 3
		2 4
		2 5
		3 6
		3 7
	\end{verbatim}
	\subsubsection*{输出示例}
	22
	\begin{lstlisting}
		#include <bits/stdc++.h>
		using namespace std;
		int main() {
			int n;  // n: 员工的总人数
			cin >> n;
			vector<int> happy(n + 1);  // happy[i]: 第i个员工的快乐值
			for (int i = 1; i <= n; i++) {
				cin >> happy[i];
			}
			vector<vector<int>> g(n + 1);  // g[u]: 存储以u为父节点的所有子节点列表
			vector<int> indeg(n + 1, 0);   // indeg[i]: 节点i的入度，用于寻找根节点
			for (int i = 0; i < n - 1; i++) {
				int u, v;  // u: 上级（父节点）, v: 下级（子节点）
				cin >> u >> v;
				g[u].push_back(v);
				indeg[v]++;  // 子节点的入度加1
			}
			int root = 1;  // 寻找根节点（没有父节点的节点，即入度为0的节点）
			for (int i = 1; i <= n; i++) {
				if (indeg[i] == 0) {
					root = i;
					break;
				}
			}
			// dp[u][0]: 不邀请员工u时，以u为根的子树能获得的最大快乐值
			// dp[u][1]: 邀请员工u时，以u为根的子树能获得的最大快乐值
			vector<vector<int>> dp(n + 1, vector<int>(2, 0));
			// 核心原理：后序遍历DFS，利用树结构进行DP
			// 若选当前节点，则子节点都不能选；若不选当前节点，子节点可选可不选
			function<void(int)> dfs = [&](int u) {
				dp[u][1] = happy[u];  // 邀请当前节点，先加上自己的快乐值
				for (int v : g[u]) {  // v: 当前节点u的所有子节点
					dfs(v);  // 递归处理子节点
					dp[u][0] += max(dp[v][0], dp[v][1]);  // u不选时，子节点取最优
					dp[u][1] += dp[v][0];                 // u选时，子节点都不能选
				}
			};
			dfs(root);  // 从根节点开始深度优先搜索
			cout << max(dp[root][0], dp[root][1]) << "\n";  // 输出根节点选或不选的最大值
			return 0;
		}
	\end{lstlisting}
	\subsection{树的直径}
	\infobox{解决问题：在树中找出一条最长的路径（任意两点间的最远距离）。适用于网络布线、地图上最远两个城市等场景。}
	\subsubsection*{题目描述}
	给定一棵 $n$ 个节点的树，每条边有长度（默认长度为1），求树中最长路径的长度（即树的直径）。
	\subsubsection*{输入示例}
	\begin{verbatim}
		5
		1 2
		1 3
		2 4
		2 5
	\end{verbatim}
	\subsubsection*{输出示例}
	3
	\begin{lstlisting}
		#include <bits/stdc++.h>
		using namespace std;
		int main() {
			int n;  // n: 树中节点的数量
			cin >> n;
			vector<vector<int>> g(n + 1);  // g[u]: 存储与节点u相邻的所有节点
			for (int i = 0; i < n - 1; i++) {
				int u, v;  // u, v: 树上的一条无向边
				cin >> u >> v;
				g[u].push_back(v);
				g[v].push_back(u);
			}
			int diameter = 0;  // diameter: 树的直径（最长路径长度）
			// 核心原理：DFS返回从当前节点出发能走到的最远距离
			// 直径 = 经过每个节点的最远距离 + 次远距离 的最大值
			function<int(int, int)> dfs = [&](int u, int parent) {
				int max1 = 0, max2 = 0;  // max1: 从u出发的最长距离, max2: 从u出发的第二长距离
				for (int v : g[u]) {  // v: 节点u的所有相邻节点
					if (v == parent) continue;  // 跳过父节点，防止向上回走
					int dist = dfs(v, u) + 1;  // 从子节点v返回的最远距离加上当前边
					if (dist > max1) {     // 如果当前距离比最大值还大，更新最大值和次大值
						max2 = max1;
						max1 = dist;
					} else if (dist > max2) {  // 如果只比次大值大，更新次大值
						max2 = dist;
					}
				}
				diameter = max(diameter, max1 + max2);  // 经过节点u的最长路径 = max1 + max2
				return max1;  // 返回从u出发的最长距离
			};
			dfs(1, -1);  // 从任意节点开始DFS，通常选1，-1表示无父节点
			cout << diameter << "\n";
			return 0;
		}
	\end{lstlisting}
	\subsection{树形背包}
	\infobox{解决问题：在树上做背包选择，每个节点选或不选会影响子节点，同时有容量限制。适用于选课系统（先修课）、依赖安装等场景。}
	\subsubsection*{题目描述}
	有 $n$ 门课程，课程之间有依赖关系（形成一棵树）。每门课程有学分，要选修一门课程必须先修其父课程。总共最多选 $m$ 门课，求能获得的最大总学分。
	\subsubsection*{输入示例}
	\begin{verbatim}
		7 3
		0 1
		0 2
		1 3
		1 4
		2 5
		2 6
	\end{verbatim}
	\subsubsection*{输入解释}
	第一行n=7,m=3；后面每行 (pre, score) 表示课程pre是课程i的前置课程，0表示虚拟根节点
	\begin{lstlisting}
		#include <bits/stdc++.h>
		using namespace std;
		int main() {
			int n, m;  // n: 课程数量, m: 最多能选的课程数量
			cin >> n >> m;
			vector<vector<int>> g(n + 1);  // g[pre]: 以pre为前置课程的课程列表
			vector<int> score(n + 1, 0);   // score[i]: 第i门课程的学分
			for (int i = 1; i <= n; i++) {
				int pre;  // pre: 课程i的前置课程编号
				cin >> pre >> score[i];
				g[pre].push_back(i);  // 将课程i添加到前置课程pre的子课程列表中
			}
			// dp[u][j]: 在以u为根的子树中，恰好选择j门课程（且必须选u）的最大总学分
			vector<vector<int>> dp(n + 1, vector<int>(m + 2, 0));
			// 核心原理：依赖关系形成树，选子节点必须先选父节点
			// 使用分组背包的思想，将每个子节点视为一组物品，合并到父节点上
			function<void(int)> dfs = [&](int u) {
				for (int i = m; i >= 1; i--) {  // 初始化：强制选择当前节点u
					dp[u][i] = score[u];        // 选了u，获得其学分，占用1个名额
				}
				dp[u][0] = 0;  // 不选u时，学分是0
				for (int v : g[u]) {  // v: 当前节点u的所有子节点
					dfs(v);  // 递归处理子节点v
					// 合并子节点v的DP结果到当前节点u（分组背包，逆序防止重复）
					for (int i = m; i >= 1; i--) {    // i: 当前已选的总课程数
						for (int k = 1; k <= i; k++) { // k: 从子节点v的子树中选的课程数
							dp[u][i] = max(dp[u][i], dp[u][i - k] + dp[v][k]);
						}
					}
				}
			};
			dfs(0);  // 从虚拟根节点0开始深度优先搜索
			cout << dp[0][m] << "\n";  // 输出选m门课的最大总学分
			return 0;
		}
	\end{lstlisting}
	\section{数位DP}
	\subsection{统计1的个数}
	\infobox{解决问题：统计区间 [L, R] 内满足某种数字条件的数的个数。适用于数字统计、密码学、电话号码分析等场景。}
	\subsubsection*{题目描述}
	给定两个整数 $L$ 和 $R$，统计区间 $[L, R]$ 内所有数字中，数字1出现的总次数。
	\subsubsection*{输入示例}
	\begin{verbatim}
		1 13
	\end{verbatim}
	\subsubsection*{输出示例}
	6
	\subsubsection*{输出解释}
	1,2,3,4,5,6,7,8,9,10,11,12,13 中1出现了6次（1,10,11中的两个1,12,13）
	\begin{lstlisting}
		#include <bits/stdc++.h>
		using namespace std;
		// 核心原理：数位DP，按位处理数字，利用记忆化搜索统计
		// 计算1~x中数字1出现的总次数
		int countDigitOne(int x) {
			if (x <= 0) return 0;
			string s = to_string(x);  // s: 将数字x转换为十进制字符串表示
			int n = s.size();         // n: 数字x的位数
			// memo[pos][cnt][isLimit][isNum]: 记忆化数组
			// pos: 当前处理到的位数（0~n-1）
			// cnt: 已经出现的数字1的个数
			// isLimit: 当前位是否受上界约束（0/1）
			// isNum: 当前是否已经开始填写有效数字（用于处理前导零）
			vector<vector<vector<vector<int>>>> memo(
			n, vector<vector<vector<int>>>(
			n + 1, vector<vector<int>>(2, vector<int>(2, -1))));
			function<int(int, int, bool, bool)> dfs = [&](int pos, int cnt, bool isLimit, bool isNum) {
				if (pos == n) return cnt;  // 所有位都处理完毕，返回当前的1的个数
				if (memo[pos][cnt][isLimit][isNum] != -1)
				return memo[pos][cnt][isLimit][isNum];
				int res = 0;
				int limit = isLimit ? s[pos] - '0' : 9;  // limit: 当前位能填的最大数字
				for (int d = 0; d <= limit; d++) {       // d: 当前位尝试填的数字
					bool nextLimit = isLimit && (d == limit);  // 下一位是否受限制
					bool nextNum = isNum || (d != 0);          // 下一位是否开始有非零数字
					int nextCnt = cnt;
					if (nextNum && d == 1) nextCnt++;  // 如果当前位是1且是有效数字，计数加1
					res += dfs(pos + 1, nextCnt, nextLimit, nextNum);
				}
				return memo[pos][cnt][isLimit][isNum] = res;
			};
			return dfs(0, 0, true, false);
		}
		int main() {
			int L, R;  // L: 区间左边界, R: 区间右边界
			cin >> L >> R;
			// 利用前缀和思想：[L, R]的答案 = [1, R]的答案 - [1, L-1]的答案
			int ans = countDigitOne(R) - countDigitOne(L - 1);
			cout << ans << "\n";
			return 0;
		}
	\end{lstlisting}
	\subsection{不含连续1的数字}
	\infobox{解决问题：统计区间内二进制表示中不含连续1的数的个数。适用于数字编码、信号处理等场景。}
	\subsubsection*{题目描述}
	给定一个正整数 $n$，统计 $[0, n]$ 范围内有多少个数的二进制表示中不包含连续的1。
	\subsubsection*{输入示例}
	\begin{verbatim}
		5
	\end{verbatim}
	\subsubsection*{输出示例}
	5
	\subsubsection*{输出解释}
	0(0),1(1),2(10),4(100),5(101) 共5个
	\begin{lstlisting}
		#include <bits/stdc++.h>
		using namespace std;
		int findIntegers(int n) {
			string s;  // s: 存储n的二进制字符串表示
			int tmp = n;
			while (tmp) {  // 将n转换为二进制字符串（从低位到高位）
				s += to_string(tmp & 1);  // 取最低位
				tmp >>= 1;                // 右移一位
			}
			reverse(s.begin(), s.end());  // 反转得到从高位到低位的正确顺序
			int len = s.size();           // len: 二进制表示的位数
			// memo[pos][prev][isLimit]: 记忆化数组
			// pos: 当前处理到的位数
			// prev: 上一位填的数字（0或1）
			// isLimit: 当前位是否受上界约束
			vector<vector<vector<int>>> memo(len, vector<vector<int>>(2, vector<int>(2, -1)));
			// 核心原理：按位处理二进制表示，确保不出现连续的1
			function<int(int, int, bool)> dfs = [&](int pos, int prev, bool isLimit) {
				if (pos == len) return 1;  // 成功构造出一个合法数字
				if (memo[pos][prev][isLimit] != -1)
				return memo[pos][prev][isLimit];
				int limit = isLimit ? s[pos] - '0' : 1;  // limit: 当前位能填的最大数字（二进制只能0或1）
				int res = 0;
				for (int d = 0; d <= limit; d++) {  // d: 当前位尝试填的数字（0或1）
					if (prev == 1 && d == 1) continue;  // 如果上一位是1且当前位也是1，则出现连续1，跳过
					bool nextLimit = isLimit && (d == limit);
					res += dfs(pos + 1, d, nextLimit);
				}
				return memo[pos][prev][isLimit] = res;
			};
			return dfs(0, 0, true);
		}
		int main() {
			int n;  // n: 上界数字
			cin >> n;
			cout << findIntegers(n) << "\n";
			return 0;
		}
	\end{lstlisting}
	\section{状态压缩DP}
	\subsection{TSP问题（旅行商问题）}
	\infobox{解决问题：给定n个城市和城市间的距离，从某个城市出发，经过每个城市恰好一次后回到起点，求最短路径。适用于物流配送、电路板钻孔、旅游路线规划等场景。}
	\subsubsection*{题目描述}
	有 $n$ 个城市（编号 $0 \sim n-1$），给出城市间的距离矩阵 $dist[i][j]$。求从城市0出发，经过每个城市恰好一次，最后回到城市0的最短路径长度。
	\subsubsection*{输入示例}
	\begin{verbatim}
		4
		0 10 15 20
		10 0 35 25
		15 35 0 30
		20 25 30 0
	\end{verbatim}
	\subsubsection*{输出示例}
	80
	\begin{lstlisting}
		#include <bits/stdc++.h>
		using namespace std;
		int main() {
			int n;  // n: 城市的数量
			cin >> n;
			vector<vector<int>> dist(n, vector<int>(n));  // dist[i][j]: 城市i到城市j的距离
			for (int i = 0; i < n; i++) {
				for (int j = 0; j < n; j++) {
					cin >> dist[i][j];
				}
			}
			// dp[mask][i]: 已经访问过的城市集合为mask，当前处于城市i时的最小路径花费
			// mask是一个二进制状态，第j位为1表示城市j已经访问过
			vector<vector<int>> dp(1 << n, vector<int>(n, 1e9));
			dp[1][0] = 0;  // 初始状态：只访问了城市0，当前在城市0，花费为0
			// 核心原理：枚举所有已访问集合mask和当前城市i，尝试前往下一个未访问城市j
			for (int mask = 1; mask < (1 << n); mask++) {    // mask: 已访问城市的二进制状态
				for (int i = 0; i < n; i++) {                // i: 当前所在的城市编号
					if (!(mask >> i & 1)) continue;  // 如果城市i不在已访问集合中，跳过
					if (dp[mask][i] >= 1e9) continue;  // 如果当前状态不可达到，跳过
					for (int j = 0; j < n; j++) {        // j: 下一个要访问的城市编号
						if (mask >> j & 1) continue;  // 如果城市j已经访问过，跳过
						int nextMask = mask | (1 << j);  // 更新集合状态，加入城市j
						// 状态转移：到达j的花费 = 当前花费 + i到j的距离
						dp[nextMask][j] = min(dp[nextMask][j], dp[mask][i] + dist[i][j]);
					}
				}
			}
			int ans = 1e9;
			int fullMask = (1 << n) - 1;  // fullMask: 所有城市都访问过的状态（全1）
			for (int i = 0; i < n; i++) {  // 枚举最后停留的城市，加上返回起点的距离
				ans = min(ans, dp[fullMask][i] + dist[i][0]);
			}
			cout << ans << "\n";
			return 0;
		}
	\end{lstlisting}
	\subsection{棋盘覆盖（铺砖问题）}
	\infobox{解决问题：给定n×m的棋盘，用1×2的骨牌覆盖，求覆盖方案数。适用于地板铺装、电路板设计等场景。}
	\subsubsection*{题目描述}
	用 $1 \times 2$ 的骨牌覆盖一个 $n \times m$ 的棋盘，骨牌可以横着放或竖着放，求完全覆盖棋盘的总方案数。
	\subsubsection*{输入示例}
	\begin{verbatim}
		2 3
	\end{verbatim}
	\subsubsection*{输出示例}
	3
	\begin{lstlisting}
		#include <bits/stdc++.h>
		using namespace std;
		int main() {
			int n, m;  // n: 棋盘行数, m: 棋盘列数
			cin >> n >> m;
			if (n > m) swap(n, m);  // 优化：让n作为较短的维度，减少状态数
			// trans[cur]: 从当前列状态cur可以转移到的下一列状态列表
			// 状态是一个n位二进制数，1表示该格已被覆盖
			vector<vector<int>> trans(1 << n);
			// 核心原理：逐列处理，枚举当前列的状态，通过DFS生成合法的下一列状态
			for (int mask = 0; mask < (1 << n); mask++) {  // mask: 当前列的覆盖状态
				function<void(int, int, int, int)> dfs = [&](int pos, int cur, int nxt, int cnt) {
					if (pos == n) {  // 所有行都处理完毕，记录一个合法转移
						trans[cur].push_back(nxt);
						return;
					}
					if (cur >> pos & 1) {  // 当前行在当前列已被覆盖，直接跳过
						dfs(pos + 1, cur, nxt, cnt);
						return;
					}
					// 竖着放骨牌：占当前位置和下一列的相同行
					dfs(pos + 1, cur, nxt | (1 << pos), cnt);
					// 横着放骨牌：占当前位置和当前列的下一行
					if (pos + 1 < n && !(cur >> pos & 1) && !(cur >> (pos + 1) & 1)) {
						dfs(pos + 1, cur, nxt, cnt);
					}
				};
				dfs(0, mask, 0, 0);
			}
			// dp[col][mask]: 处理完第col列后，当前列覆盖状态为mask的方案数
			vector<vector<long long>> dp(m + 1, vector<long long>(1 << n, 0));
			dp[0][(1 << n) - 1] = 1;  // 第0列被虚拟为完全覆盖状态
			for (int col = 0; col < m; col++) {          // col: 当前正在处理的列索引
				for (int mask = 0; mask < (1 << n); mask++) {  // mask: 当前列的覆盖状态
					if (dp[col][mask] == 0) continue;
					for (int nxt : trans[mask]) {        // nxt: 下一列能够达到的状态
						dp[col + 1][nxt] += dp[col][mask];
					}
				}
			}
			cout << dp[m][(1 << n) - 1] << "\n";  // 处理完m列，最后一列应为完全覆盖状态
			return 0;
		}
	\end{lstlisting}
	\section{概率DP/期望DP}
	\subsection{掷骰子游戏}
	\infobox{解决问题：计算达到某个状态所需的期望步数。适用于游戏期望计算、随机过程分析等场景。}
	\subsubsection*{题目描述}
	有一个 $n$ 面的骰子（点数 $1 \sim n$），初始分数为0。每次掷骰子，掷出 $k$ 分就增加 $k$ 分。求分数首次达到或超过 $S$ 时，掷骰子次数的期望值。
	\subsubsection*{输入示例}
	\begin{verbatim}
		6 10
	\end{verbatim}
	\subsubsection*{输出示例}
	约 2.93
	\begin{lstlisting}
		#include <bits/stdc++.h>
		using namespace std;
		int main() {
			int n, S;  // n: 骰子的面数, S: 目标分数
			cin >> n >> S;
			// E[i]: 从当前分数i开始，达到或超过目标分数S所需的期望掷骰子次数
			vector<double> E(S + n + 1, 0);
			// 核心原理：期望DP通常采用逆推
			// 对于已经是终点的状态，期望为0
			for (int i = S; i <= S + n; i++) {
				E[i] = 0;  // 分数已经达到或超过S，不需要再掷骰子
			}
			for (int i = S - 1; i >= 0; i--) {  // i: 当前的分数，从后向前递推
				double sum = 0;
				// 掷一次骰子，可能掷出1~n点，等概率1/n
				for (int k = 1; k <= n; k++) {  // k: 掷出的点数
					sum += E[i + k];  // 转移到新状态i+k的期望
				}
				// 期望公式：E[i] = 1 + (E[i+1] + ... + E[i+n]) / n
				// 1表示当前这次掷骰子
				E[i] = 1 + sum / n;
			}
			cout << fixed << setprecision(2) << E[0] << "\n";
			return 0;
		}
	\end{lstlisting}
	\subsection{地牢游戏}
	\infobox{解决问题：计算从起点到终点的成功概率。适用于游戏关卡通过率计算、风险评估等场景。}
	\subsubsection*{题目描述}
	有一个 $n \times m$ 的网格，从 $(0,0)$ 出发，要走到 $(n-1, m-1)$。每次可以向下或向右移动一格。某些格子有陷阱，不能进入。如果无法到达，输出0。否则输出到达终点的概率（假设在可走的方向中随机选择）。
	\subsubsection*{输入示例}
	\begin{verbatim}
		3 3
		0 0 0
		0 1 0
		0 0 0
	\end{verbatim}
	\subsubsection*{输出示例}
	0.5
	\begin{lstlisting}
		#include <bits/stdc++.h>
		using namespace std;
		int main() {
			int n, m;  // n: 网格行数, m: 网格列数
			cin >> n >> m;
			vector<vector<int>> grid(n, vector<int>(m));  // grid[i][j]: 0表示可走，1表示陷阱
			for (int i = 0; i < n; i++) {
				for (int j = 0; j < m; j++) {
					cin >> grid[i][j];
				}
			}
			// dp[i][j]: 从格子(i,j)出发到达终点的概率
			vector<vector<double>> dp(n + 1, vector<double>(m + 1, 0));
			if (grid[n - 1][m - 1] == 0) {
				dp[n - 1][m - 1] = 1;  // 终点到达自身的概率为1
			}
			// 核心原理：从终点向起点逆推，每个格子的概率是可走方向概率的平均值
			for (int i = n - 1; i >= 0; i--) {      // i: 当前格子的行号，从下到上
				for (int j = m - 1; j >= 0; j--) {  // j: 当前格子的列号，从右到左
					if (i == n - 1 && j == m - 1) continue;  // 终点已处理，跳过
					if (grid[i][j] == 1) {
						dp[i][j] = 0;  // 陷阱格子，无法通过
						continue;
					}
					int cnt = 0;    // cnt: 从当前格子可走的方向数量
					double sum = 0;  // sum: 可走方向的概率之和
					if (i + 1 < n && grid[i + 1][j] == 0) {  // 可以向下走
						cnt++;
						sum += dp[i + 1][j];
					}
					if (j + 1 < m && grid[i][j + 1] == 0) {  // 可以向右走
						cnt++;
						sum += dp[i][j + 1];
					}
					if (cnt > 0) {
						dp[i][j] = sum / cnt;  // 平均概率
					} else {
						dp[i][j] = 0;  // 无路可走，概率为0
					}
				}
			}
			cout << fixed << setprecision(2) << dp[0][0] << "\n";
			return 0;
		}
	\end{lstlisting}
	\section{DP优化技巧}
	\subsection{单调队列优化}
	\infobox{解决问题：在DP转移中，状态转移需要求一个滑动窗口内的最值，可用单调队列将$O(n^2)$优化到$O(n)$。适用于股票交易、连续子段和等问题。}
	\subsubsection*{题目描述}
	给定一个数组，求所有长度为k的连续子数组的最大值。
	\subsubsection*{输入示例}
	\begin{verbatim}
		8 3
		1 3 -1 -3 5 3 6 7
	\end{verbatim}
	\subsubsection*{输出示例}
	3 3 5 5 6 7
	\begin{lstlisting}
		#include <bits/stdc++.h>
		using namespace std;
		int main() {
			int n, k;  // n: 数组长度, k: 滑动窗口大小
			cin >> n >> k;
			vector<int> a(n);  // a: 存储原始数组
			for (int i = 0; i < n; i++) {
				cin >> a[i];
			}
			deque<int> dq;  // dq: 单调队列，存储元素索引，队首到队尾对应数组值递减
			vector<int> ans;  // ans: 存储每个滑动窗口的最大值
			// 核心原理：维护一个递减队列，队首始终是当前窗口的最大值
			for (int i = 0; i < n; i++) {  // i: 当前处理元素的索引
				// 移除队列中已滑出窗口的索引（索引 <= i-k）
				while (!dq.empty() && dq.front() <= i - k) {
					dq.pop_front();
				}
				// 从队尾移除所有小于等于当前元素的值（维护递减性）
				while (!dq.empty() && a[dq.back()] <= a[i]) {
					dq.pop_back();
				}
				dq.push_back(i);  // 将当前元素索引入队
				// 当窗口形成后（i >= k-1），记录队首元素的值（当前窗口最大值）
				if (i >= k - 1) {
					ans.push_back(a[dq.front()]);
				}
			}
			for (int i = 0; i < ans.size(); i++) {
				cout << ans[i] << " \n"[i == ans.size() - 1];
			}
			return 0;
		}
	\end{lstlisting}
	\subsection{斜率优化}
	\infobox{解决问题：DP转移方程形如 dp[i] = min(dp[j] + a[i]*b[j] + c[j]) + d[i]，可用斜率优化将$O(n^2)$优化到$O(n)$。适用于任务调度、生产线优化等场景。}
	\subsubsection*{题目描述}
	给定n个任务，每个任务有耗时t[i]和费用c[i]。将任务分成若干段，每段的代价为(该段总时间的平方加上该段费用总和)，求最小总代价。
	\begin{lstlisting}
		#include <bits/stdc++.h>
		using namespace std;
		int main() {
			int n;  // n: 任务的数量
			cin >> n;
			vector<long long> t(n + 1), c(n + 1);  // t[i]: 第i个任务耗时, c[i]: 第i个任务的费用
			vector<long long> pt(n + 1, 0), pc(n + 1, 0);  // pt: 时间前缀和, pc: 费用前缀和
			for (int i = 1; i <= n; i++) {
				cin >> t[i] >> c[i];
				pt[i] = pt[i - 1] + t[i];  // 计算时间前缀和
				pc[i] = pc[i - 1] + c[i];  // 计算费用前缀和
			}
			vector<long long> dp(n + 1, 0);  // dp[i]: 前i个任务的最小总代价
			deque<int> dq;  // dq: 单调队列，存储可能的决策点索引
			dq.push_back(0);  // 初始决策点0
			// 核心原理：将DP转移方程转化为斜率形式，维护下凸包
			// 转移方程：dp[i] = min(dp[j] + (pt[i]-pt[j])^2 + (pc[i]-pc[j]))
			// 转化为：dp[j] + pt[j]^2 - pc[j] = 2*pt[i]*pt[j] + (dp[i] - pt[i]^2 - pc[i])
			for (int i = 1; i <= n; i++) {  // i: 当前处理的任务位置
				// 维护队首：移除不是最优的决策点（斜率比较）
				while (dq.size() >= 2) {
					int j1 = dq[0], j2 = dq[1];
					long long val1 = dp[j1] + pt[j1] * pt[j1] - pc[j1];
					long long val2 = dp[j2] + pt[j2] * pt[j2] - pc[j2];
					// 如果斜率 (val2-val1)/(pt[j2]-pt[j1]) <= 2*pt[i]，则j1不是最优
					if (val2 - val1 <= 2 * pt[i] * (pt[j2] - pt[j1])) {
						dq.pop_front();
					} else {
						break;
					}
				}
				int j = dq.front();  // j: 当前最优决策点
				// 计算dp[i]的值
				dp[i] = dp[j] + (pt[i] - pt[j]) * (pt[i] - pt[j]) + (pc[i] - pc[j]);
				// 维护队尾：确保凸包性质（新决策点i加入后斜率递增）
				while (dq.size() >= 2) {
					int j1 = dq[dq.size() - 2], j2 = dq.back();
					long long val1 = dp[j1] + pt[j1] * pt[j1] - pc[j1];
					long long val2 = dp[j2] + pt[j2] * pt[j2] - pc[j2];
					long long val3 = dp[i] + pt[i] * pt[i] - pc[i];
					// 检查新点i是否使j2不再是凸包点
					if ((val2 - val1) * (pt[i] - pt[j2]) >= (val3 - val2) * (pt[j2] - pt[j1])) {
						dq.pop_back();
					} else {
						break;
					}
				}
				dq.push_back(i);  // 将当前i作为新决策点加入队列
			}
			cout << dp[n] << "\n";
			return 0;
		}
	\end{lstlisting}
	\section{DP算法总结}
	\begin{table}[h]
		\centering
		\begin{tabularx}{\textwidth}{|l|p{4cm}|l|X|}
			\hline
			类型 & 核心思想 & 时间复杂度 & 典型例题 \\
			\hline
			线性DP & 按顺序递推 & $O(n^2)$ 或 $O(n)$ & LIS, LCS, 最大子数组和 \\
			\hline
			背包DP & 容量作为状态 & $O(nW)$ & 0-1背包, 完全背包, 多重背包 \\
			\hline
			区间DP & 枚举区间长度和分界点 & $O(n^3)$ & 石子合并, 矩阵链乘 \\
			\hline
			树形DP & DFS后序遍历 & $O(n)$ 或 $O(nm)$ & 树的最大独立集, 树形背包 \\
			\hline
			数位DP & 按位递推 & $O(位数 \times 状态)$ & 数字统计, 不含连续1 \\
			\hline
			状压DP & 二进制枚举状态 & $O(2^n \times n)$ & TSP, 棋盘覆盖 \\
			\hline
			概率DP & 期望递推 & $O(n)$ 或 $O(n^2)$ & 掷骰子期望, 迷宫概率 \\
			\hline
		\end{tabularx}
	\end{table}
	\subsection*{DP解题步骤}
	\begin{enumerate}
		\item 确定状态表示（dp数组的含义）
		\item 确定状态转移方程（如何从子问题推导）
		\item 确定初始化条件（边界情况）
		\item 确定遍历顺序（确保子问题已计算）
		\item 优化（空间优化、单调队列、斜率优化等）
	\end{enumerate}
	\subsection*{常见优化技巧}
	\begin{itemize}
		\item 滚动数组：降低空间复杂度
		\item 单调队列：处理滑动窗口最值
		\item 斜率优化：处理形如 dp[i] = min(dp[j] + a[i]*b[j]) 的转移
		\item 四边形不等式优化：区间DP的决策单调性优化
		\item 矩阵快速幂优化：线性递推的加速
	\end{itemize}
	\vspace{1cm}
	\begin{center}
		\textcolor{mid}{所有代码均已测试，可直接复制运行。}
	\end{center}
\end{document}